\section{Methodology}
\label{sec:methodology}

\subsection{Geometry-Conditioned Yoga Pose Understanding}
\label{sec:method_overview}

We propose \textbf{YogaLLaVA V3}, the paper-facing name used for the
PoseVLM V3 implementation and its selected checkpoint. YogaLLaVA V3 is
a geometry-conditioned vision--language model for fine-grained yoga-pose
understanding. Given an RGB image $I$, a natural-language question $q$,
and an ordered set of structured body measurements $\mathcal{G}$, the
model generates a natural-language answer $y$ that may contain both
qualitative explanations and continuous angular values. We formulate
the conditional generation problem as
\begin{equation}
    p(y \mid I,q,\mathcal{G}).
    \label{eq:problem_formulation}
\end{equation}

The structured measurements are produced by the upstream SAM 3D Body
reconstruction and metric-extraction pipeline described in
Sec.~\ref{sec:dataset}. YogaLLaVA V3 therefore performs
language-conditioned reasoning over visual information and supplied
geometry. It does not independently reconstruct the complete geometry
vector from RGB pixels.

The model uses a fixed vocabulary of $K=14$ degree-valued descriptors:
left and right knee flexion, left and right elbow flexion, left and
right hip opening, left and right shoulder elevation, trunk tilt,
lateral trunk bend, trunk twist, neck alignment, knee symmetry, and
elbow symmetry. These descriptors form a degree-based subset of the
broader annotation taxonomy introduced in Sec.~\ref{sec:dataset}.
Scale-dependent quantities such as stance width and body-center offset
are not supplied to the model. Hip symmetry is also not represented by
a dedicated input descriptor, although three hip-symmetry targets occur
in the held-out test set.

The geometry input is represented as
\begin{equation}
    \mathcal{G}
    =
    \left\{
        \left(\theta_k,m_k\right)
    \right\}_{k=1}^{K},
    \qquad
    m_k\in\{0,1\},
    \label{eq:geometry_input}
\end{equation}
where $\theta_k$ is the value of the $k$-th descriptor in degrees and
$m_k$ indicates whether that measurement is available. The validity
indicator is modeled explicitly because an unavailable measurement is
semantically different from a valid value close to zero.

YogaLLaVA V3 is built upon LLaVA-1.5-7B and introduces two mechanisms
for continuous geometry reasoning. A geometry encoder maps the ordered
structured measurements into token representations compatible with the
language-model hidden space, while a dedicated angle-regression head
predicts continuous values at special numeric positions in the answer.
The resulting architecture integrates visual features, structured
geometry, and language context in a unified autoregressive model.

\subsection{Multimodal Geometry Representation}
\label{sec:geometry_representation}

\paragraph{Visual and linguistic representations.}
The input image is processed by the pretrained CLIP ViT-L/14-336 vision
encoder of LLaVA. The resulting patch-level visual features are mapped
into the language-model hidden space through the frozen multimodal
projector. The instruction and question are tokenized using the LLaVA
tokenizer, and the standard LLaVA multimodal mechanism integrates the
projected visual features with the textual sequence.

To incorporate structured body geometry, the prompt contains exactly
$K$ special \texttt{<GEO>} placeholders arranged in canonical metric
order. Their ordinary token embeddings are replaced by the corresponding
geometry representations before the sequence is processed by the causal
language decoder. This places visual, linguistic, and structured numeric
information in a shared embedding space.

\paragraph{Continuous scalar encoding.}
Each available degree-valued descriptor is normalized to $[0,1]$:
\begin{equation}
    v_k = \frac{\theta_k}{180}.
    \label{eq:angle_normalization}
\end{equation}
To obtain a richer representation than a direct scalar projection, we
encode $v_k$ using $B=16$ fixed logarithmically spaced Fourier frequency
bands. The scalar feature map is
\begin{equation}
\begin{split}
    \boldsymbol{\phi}(v_k)
    =
    \big[
        &\sin(\omega_1v_k),\ldots,\sin(\omega_Bv_k),\\
        &\cos(\omega_1v_k),\ldots,\cos(\omega_Bv_k),
        v_k
    \big].
\end{split}
\label{eq:fourier_features}
\end{equation}
This construction yields $2B+1=33$ features. A two-layer multilayer
perceptron with hidden dimension 256 projects the Fourier representation
into the $d=4096$ dimensional language-model space:
\begin{equation}
    \mathbf{z}_k
    =
    f_{\mathrm{geo}}
    \left(
        \boldsymbol{\phi}(v_k)
    \right),
    \qquad
    \mathbf{z}_k\in\mathbb{R}^{d}.
    \label{eq:scalar_projection}
\end{equation}

\paragraph{Metric identity and missing geometry.}
A scalar value is not meaningful without its anatomical identity. Each
descriptor therefore has a learned metric-identity embedding
$\mathbf{e}^{\mathrm{id}}_k\in\mathbb{R}^{d}$. The geometry token for
descriptor $k$ is
\begin{equation}
    \mathbf{g}_k =
    \begin{cases}
    \operatorname{LN}\!\left(
        \mathbf{e}^{\mathrm{id}}_k + \mathbf{z}_k
    \right), & m_k=1, \\[5pt]
    \operatorname{LN}\!\left(
        \mathbf{e}^{\mathrm{id}}_k +
        \mathbf{e}^{\mathrm{miss}}
    \right), & m_k=0,
    \end{cases}
    \label{eq:geometry_token}
\end{equation}
where $\mathbf{e}^{\mathrm{miss}}\in\mathbb{R}^{d}$ is a learned
missing-value embedding and $\operatorname{LN}$ denotes layer
normalization. The complete geometry sequence is
\begin{equation}
    \mathbf{G}
    =
    [\mathbf{g}_1,\ldots,\mathbf{g}_K]
    \in\mathbb{R}^{K\times d}.
    \label{eq:geometry_sequence}
\end{equation}
These vectors replace the corresponding \texttt{<GEO>} placeholder
embeddings, allowing the decoder to jointly attend to image features,
geometry tokens, and language context.

\subsection{Angle-Aware Language Generation}
\label{sec:angle_generation}

Conventional language models express numbers as sequences of discrete
tokens. To separate linguistic structure from continuous-value
prediction, YogaLLaVA V3 introduces the special token
\texttt{<ANGLE>}. During training, each degree value in a
geometry-grounded answer is replaced by \texttt{<ANGLE>}, while the
corresponding normalized target is stored separately and aligned with
that token position. For example,
\begin{quote}
\small
\textit{The left knee flexion angle is $140.2^\circ$.}
\end{quote}
is represented internally as
\begin{quote}
\small
\texttt{The left knee flexion angle is <ANGLE>.}
\end{quote}

The language model learns where a numeric value belongs by predicting
\texttt{<ANGLE>}; the regression branch predicts the continuous number
associated with that position. Let $p_j$ denote the sequence position of
the $j$-th \texttt{<ANGLE>} token. Because the decoder is causal, the
hidden state $\mathbf{h}_{p_j-1}$ is the state that produces the token at
$p_j$. The normalized angle prediction is
\begin{equation}
    \widehat{v}_j
    =
    \sigma\!\left(
        f_{\mathrm{reg}}
        \left(
            \mathbf{h}_{p_j-1}
        \right)
    \right),
    \label{eq:angle_regression}
\end{equation}
where $f_{\mathrm{reg}}$ consists of layer normalization, a
$4096\!\rightarrow\!512$ linear layer, GELU activation, and a
$512\!\rightarrow\!1$ linear layer. The sigmoid constrains the output to
$[0,1]$, and the value in degrees is recovered as
\begin{equation}
    \widehat{\theta}_j
    =
    180\widehat{v}_j.
    \label{eq:angle_denormalization}
\end{equation}

Reading $\mathbf{h}_{p_j-1}$ ensures that the current scalar is predicted
from the preceding causal context rather than from the value-conditioned
embedding placed at the current \texttt{<ANGLE>} position.

\paragraph{Continuous-value feedback.}
After an angle slot is emitted, its scalar value is embedded back into
the sequence so that subsequent words and later angle slots can
condition on it. A second scalar encoder, independent of
$f_{\mathrm{geo}}$, defines
\begin{equation}
    \mathbf{a}(u_j)
    =
    \mathbf{e}^{\mathrm{angle}}
    +
    f_{\mathrm{slot}}
    \left(
        \boldsymbol{\phi}(u_j)
    \right),
    \label{eq:angle_feedback}
\end{equation}
where $\mathbf{e}^{\mathrm{angle}}$ is a learned base embedding. During
teacher-forced training, $u_j$ is the reference normalized value.
During autoregressive inference, $u_j=\widehat{v}_j$. Thus, later output
tokens condition on the reference value during training and on the
model's own prediction during generation.

\subsection{Joint Language and Geometry Learning}
\label{sec:training_objective}

YogaLLaVA V3 is trained jointly for answer generation and continuous
angle prediction. Language supervision is applied only to the assistant
answer span; the image placeholder, geometry placeholders, instruction,
question, and padding positions are excluded from the language loss.

Let $\mathcal{T}$ denote the supervised assistant-token positions. The
causal language-modeling loss is
\begin{equation}
    \mathcal{L}_{\mathrm{LM}}
    =
    -\frac{1}{|\mathcal{T}|}
    \sum_{t\in\mathcal{T}}
    \log
    p
    \left(
        y_t
        \mid
        y_{<t},I,q,\mathcal{G}
    \right).
    \label{eq:language_loss}
\end{equation}
For a batch containing $N$ supervised angle slots, continuous prediction
is optimized in normalized angle space using
\begin{equation}
    \mathcal{L}_{\mathrm{ang}}
    =
    \frac{1}{N}
    \sum_{j=1}^{N}
    \operatorname{SmoothL1}_{\beta}
    \left(
        \widehat{v}_j-v_j
    \right).
    \label{eq:regression_loss}
\end{equation}
The complete objective is
\begin{equation}
    \mathcal{L}
    =
    \mathcal{L}_{\mathrm{LM}}
    +
    \lambda_{\mathrm{ang}}
    \mathcal{L}_{\mathrm{ang}},
    \label{eq:total_loss}
\end{equation}
with $\lambda_{\mathrm{ang}}=0.75$ and $\beta=0.05$, corresponding to a
$9^\circ$ transition point. The language term is implemented as causal
cross-entropy with label smoothing. Domain-knowledge records containing
no \texttt{<ANGLE>} slots contribute only
$\mathcal{L}_{\mathrm{LM}}$.

Full geometry dropout is applied to 15\% of geometry-grounded training
records. For a selected record, all geometry measurements are marked
unavailable and represented through the learned missing-value embedding.

The CLIP vision encoder and multimodal projector remain frozen. The
Vicuna language decoder is adapted through LoRA with rank 32, scaling
parameter 64, and dropout 0.10. LoRA is applied to the
$q$, $k$, $v$, and $o$ attention projections and the gate, up, and down
feed-forward projections. The geometry encoder, angle-feedback encoder,
learned angle base embedding, regression head, and LoRA parameters are
trained jointly. The input embedding and language-model output
projection are also trainable after adding \texttt{<GEO>} and
\texttt{<ANGLE>} to the vocabulary.

\subsection{Autoregressive Inference}
\label{sec:inference}

At inference time, the upstream reconstruction pipeline supplies the
ordered geometry descriptors, while the image is processed by the
LLaVA visual stream. The image, geometry tokens, instruction, and
question are provided without a reference answer.

Generation proceeds one token at a time. When an ordinary token is
selected, decoding continues normally. When the model emits
\texttt{<ANGLE>}, the hidden state that produced the token is passed
through the regression head to obtain a continuous prediction. The
prediction is converted to degrees, encoded through
Eq.~\eqref{eq:angle_feedback}, and inserted as the embedding of the
emitted \texttt{<ANGLE>} token before generation continues. After
decoding, the \texttt{<ANGLE>} placeholders are replaced in sequence
order by their predicted degree values. Token probabilities and
regression values are not averaged; they interact through this
placeholder-and-regression mechanism.

\section{Experimental Setup}
\label{sec:experimental_setup}

\subsection{Dataset Splits}
\label{sec:dataset_splits}

We conduct experiments on the combined domain-knowledge and
geometry-grounded QA dataset introduced in Sec.~\ref{sec:dataset}. The
dataset contains 10,090 records: 6,000 geometry-grounded examples and
4,090 domain-knowledge examples. The fixed splits are summarized in
Table~\ref{tab:dataset_splits}.

\begin{table}[t]
    \centering
    \caption{Composition of the YogaLLaVA V3 instruction-tuning dataset.}
    \label{tab:dataset_splits}
    \resizebox{\columnwidth}{!}{
    \begin{tabular}{lrrrr}
        \toprule
        \textbf{Split} &
        \textbf{Geometry QA} &
        \textbf{Domain QA} &
        \textbf{Total} &
        \textbf{Angle Slots} \\
        \midrule
        Training   & 5,396 & 3,690 & 9,086 & 9,163 \\
        Validation &   294 &   200 &   494 &   493 \\
        Test       &   310 &   200 &   510 &   521 \\
        \midrule
        Total      & 6,000 & 4,090 & 10,090 & 10,177 \\
        \bottomrule
    \end{tabular}}
\end{table}

The geometry-grounded subset is constructed from 3,000 selected images,
with two QA records associated with each image. It contains 3,000
direct-angle, 905 comparison, 900 regional-summary, 600 misconception,
and 595 instructor-observation records.

The test images are not present in training, but pose categories are
shared across the splits. The evaluation therefore measures
generalization to unseen images within known yoga-pose categories, not
generalization to entirely unseen poses.

\subsection{Model Configuration}
\label{sec:model_configuration}

YogaLLaVA V3 is initialized from
\texttt{llava-hf/llava-1.5-7b-hf}. The vision component uses CLIP
ViT-L/14-336 at $336\times336$ resolution, yielding 576 selected patch
tokens with feature width 1,024. The multimodal projector maps these
features to the 4,096-dimensional language-model hidden space. The
Vicuna/LLaMA decoder contains 32 transformer layers and 32 attention
heads. The maximum tokenized sequence length is 1,280.

The scalar encoders use 16 Fourier bands, a 33-dimensional Fourier
representation, a hidden projection of 256, and an output projection of
4,096. The angle-regression head uses the architecture
$4096\!\rightarrow\!512\!\rightarrow\!1$ with layer normalization,
GELU, and sigmoid activation. The vision encoder and multimodal
projector are frozen, while the LoRA adapters and custom geometry and
regression modules are trained.

\subsection{Training Protocol}
\label{sec:training_protocol}

YogaLLaVA V3 is trained for 12 epochs with AdamW, weight decay 0.01,
and gradient clipping at 1.0. The per-step batch size is four, and
gradients are accumulated over eight minibatches, producing an effective
batch size of 32. We use a cosine learning-rate schedule with linear
warm-up over the first 3\% of optimizer updates.

Three parameter groups use separate learning rates:
\begin{itemize}
    \item $1.5\times10^{-4}$ for the LoRA adapters;
    \item $2\times10^{-5}$ for the input-embedding and output-projection
    matrices; and
    \item $1\times10^{-4}$ for the geometry encoder, scalar-value
    encoders, angle base embedding, and regression head.
\end{itemize}

The pretrained base model is loaded in bfloat16 precision, while the
custom numerical modules are maintained in float32. All experiments are
conducted using PyTorch on a single NVIDIA RTX A6000 GPU. Python, NumPy,
PyTorch, and CUDA random seeds are fixed before training.

Validation is performed after each epoch. The selected checkpoint is
the epoch-10 checkpoint, which achieves the lowest teacher-forced
validation angle RMSE of $19.8274^\circ$, with validation MAE
$12.6342^\circ$ and $R^2=0.8739$. The checkpoint-selection criterion
therefore optimizes teacher-forced numeric validation accuracy rather
than free-running language quality.

\subsection{Evaluation Protocol}
\label{sec:evaluation_protocol}

\paragraph{Teacher-forced angle-slot regression.}
The principal numeric test evaluates teacher-forced angle-slot
regression. The full reference answer sequence is supplied, including
the correct number, positions, and ordering of all
\texttt{<ANGLE>} tokens. The regression head predicts the numeric value
at each known slot.

This protocol evaluates continuous prediction given the reference
answer structure. It does not test whether free-running generation emits
an angle slot when required, produces the correct number of slots,
associates each slot with the intended anatomical metric, preserves
multi-metric ordering, generates coherent non-repetitive prose,
maintains consistency between numeric predictions and qualitative
conclusions, or terminates correctly.

For $N$ supervised slots with reference values $\theta_i$ and
predictions $\widehat{\theta}_i$, we report
\begin{equation}
    \operatorname{MAE}
    =
    \frac{1}{N}
    \sum_{i=1}^{N}
    \left|
        \widehat{\theta}_i-\theta_i
    \right|,
    \label{eq:evaluation_mae}
\end{equation}
\begin{equation}
    \operatorname{RMSE}
    =
    \sqrt{
        \frac{1}{N}
        \sum_{i=1}^{N}
        \left(
            \widehat{\theta}_i-\theta_i
        \right)^2
    },
    \label{eq:evaluation_rmse}
\end{equation}
and
\begin{equation}
    R^2
    =
    1-
    \frac{
        \sum_{i=1}^{N}
        \left(
            \theta_i-\widehat{\theta}_i
        \right)^2
    }{
        \sum_{i=1}^{N}
        \left(
            \theta_i-\overline{\theta}
        \right)^2
    }.
    \label{eq:evaluation_r2}
\end{equation}
We additionally report Pearson correlation and the percentages of
predictions within $5^\circ$, $10^\circ$, and $15^\circ$ of the
references.

For 518 of the 521 test slots, the exact requested measurement is
already present in the structured geometry input. This evaluation
therefore primarily measures question-conditioned descriptor
identification, routing, selection, numeric reproduction, and
contextual expression of supplied geometry. It is not an evaluation of
independent image-only physical angle recovery.

The numeric references are reconstruction-derived pseudo-measurements
from SAM 3D Body. They are not motion-capture, manually verified
biomechanical, clinical, diagnostic, or rehabilitation ground truth.

\paragraph{Autoregressive generation comparison.}
We also report free-running generation on the 310 geometry-grounded test
records. This comparison uses the selected PoseVLM V3 checkpoint
evaluated through the later comparison pipeline and explanatory prompt;
the generation metrics are separate from the original teacher-forced
test artifact. YogaLLaVA V3 receives image, question, and structured
geometry, while zero-shot LLaVA-1.5-7B and Qwen2.5-VL-7B-Instruct
receive image and question only. The baseline questions identify the
requested target metrics and expected order, but do not provide their
numeric geometry values.

Because numeric outputs can be aligned only when the generated and
reference slot counts match, we report model-specific slot-count
coverage and a common-success comparison over records on which all three
models produce the correct number of numeric slots. We do not report
modality ablations, pose-held-out evaluation, or independently annotated
biomechanical evaluation.

\section{Results and Analysis}
\label{sec:results}

\subsection{Overall Teacher-Forced Angle-Slot Regression}
\label{sec:overall_results}

Table~\ref{tab:overall_results} reports teacher-forced angle-slot
regression over 521 known reference slots from the 310
geometry-grounded test records.

\begin{table}[t]
    \centering
    \caption{Teacher-forced angle-slot regression on the held-out test
    set. Errors are measured against SAM 3D Body-derived geometric
    pseudo-measurements.}
    \label{tab:overall_results}
    \begin{tabular}{lr}
        \toprule
        \textbf{Metric} & \textbf{YogaLLaVA V3} \\
        \midrule
        Number of angle slots & 521 \\
        MSE & $476.28~\mathrm{deg}^2$ \\
        RMSE $\downarrow$ & $21.82^\circ$ \\
        MAE $\downarrow$ & $13.98^\circ$ \\
        $R^2$ $\uparrow$ & 0.832 \\
        Pearson correlation $\uparrow$ & 0.913 \\
        Within $5^\circ$ $\uparrow$ & 34.36\% \\
        Within $10^\circ$ $\uparrow$ & 57.01\% \\
        Within $15^\circ$ $\uparrow$ & 70.63\% \\
        \bottomrule
    \end{tabular}
\end{table}

At known reference slot positions, YogaLLaVA V3 differs from the
reconstruction-derived values by $13.98^\circ$ on average and achieves
an RMSE of $21.82^\circ$. Approximately 57.01\% of the predictions lie
within $10^\circ$ of the corresponding references. These values are
valid held-out teacher-forced regression results, but they are not
complete free-running inference metrics.

Because 518 of 521 targets are already available in the structured
geometry input, the result primarily supports the model's ability to
identify, route, select, reproduce, and linguistically express supplied
geometry. It does not demonstrate independent recovery of physical
joint angles from RGB pixels.

\subsection{Performance Across Question Types}
\label{sec:path_results}

\begin{table}[t]
    \centering
    \caption{Teacher-forced angle-slot regression grouped by
    geometry-question type.}
    \label{tab:path_results}
    \resizebox{\columnwidth}{!}{
    \begin{tabular}{lrrrrr}
        \toprule
        \textbf{Question Type} &
        \textbf{Slots} &
        \textbf{RMSE $\downarrow$} &
        \textbf{MAE $\downarrow$} &
        \textbf{$R^2$ $\uparrow$} &
        \textbf{Within $10^\circ$ $\uparrow$} \\
        \midrule
        Comparison & 106 & $21.53^\circ$ & $12.95^\circ$ & 0.791 & 62.26\% \\
        Direct angle & 155 & $22.65^\circ$ & $15.22^\circ$ & 0.821 & 51.61\% \\
        Instructor observation & 31 & $28.69^\circ$ & $21.08^\circ$ & 0.054 & 25.81\% \\
        Misconception & 32 & $17.10^\circ$ & $11.28^\circ$ & 0.877 & 68.75\% \\
        Region summary & 197 & $20.72^\circ$ & $12.88^\circ$ & 0.847 & 61.42\% \\
        \bottomrule
    \end{tabular}}
\end{table}

Misconception questions have the lowest MAE and the highest
within-$10^\circ$ rate, although the subset contains only 32 slots and
should be interpreted cautiously. Regional summaries also perform
relatively well despite requiring several measurements to be expressed
within a single answer.

Instructor-observation questions are the weakest category: they have
the largest RMSE and MAE, the lowest $R^2$, and the lowest
within-$10^\circ$ accuracy. One possible explanation is that indirect
or teacher-style wording makes the requested descriptor less explicit,
although the current experiment does not isolate the cause.

\subsection{Performance Across Geometric Descriptors}
\label{sec:metric_results}

\begin{table}[t]
    \centering
    \caption{Teacher-forced angle-slot regression grouped by target
    descriptor.}
    \label{tab:metric_results}
    \resizebox{\columnwidth}{!}{
    \begin{tabular}{lrrr}
        \toprule
        \textbf{Descriptor} &
        \textbf{Slots} &
        \textbf{RMSE $\downarrow$} &
        \textbf{MAE $\downarrow$} \\
        \midrule
        Knee symmetry & 4 & $3.23^\circ$ & $2.68^\circ$ \\
        Elbow symmetry & 4 & $6.90^\circ$ & $5.53^\circ$ \\
        Hip symmetry & 3 & $10.18^\circ$ & $7.82^\circ$ \\
        Trunk twist & 33 & $11.02^\circ$ & $6.84^\circ$ \\
        Lateral trunk bend & 38 & $14.92^\circ$ & $9.24^\circ$ \\
        Right shoulder elevation & 34 & $14.82^\circ$ & $9.89^\circ$ \\
        Trunk tilt & 34 & $15.48^\circ$ & $9.50^\circ$ \\
        Neck alignment & 37 & $17.34^\circ$ & $12.75^\circ$ \\
        Left hip opening & 48 & $18.40^\circ$ & $13.28^\circ$ \\
        Right knee flexion & 51 & $22.55^\circ$ & $12.09^\circ$ \\
        Left knee flexion & 47 & $23.66^\circ$ & $15.56^\circ$ \\
        Right elbow flexion & 50 & $24.80^\circ$ & $18.39^\circ$ \\
        Right hip opening & 48 & $25.41^\circ$ & $17.36^\circ$ \\
        Left elbow flexion & 47 & $27.77^\circ$ & $18.46^\circ$ \\
        Left shoulder elevation & 43 & $31.22^\circ$ & $21.23^\circ$ \\
        \bottomrule
    \end{tabular}}
\end{table}

The symmetry rows contain only three or four slots and do not support
broad conclusions. Among descriptors with more substantial test
support, trunk twist has the lowest MAE. Trunk tilt, lateral trunk bend,
and right shoulder elevation also show comparatively low errors.

Left shoulder elevation is the most difficult descriptor, with the
largest RMSE and MAE. Left and right elbow flexion and right hip opening
are also relatively difficult. Differences may reflect descriptor
distributions, reconstruction noise, template frequency, or
left--right ambiguity, but these explanations remain hypotheses rather
than established causes.

Hip symmetry is structurally different from the other rows because it
is not represented by a dedicated token in the 14-descriptor input
vocabulary. Its three targets therefore cannot be directly routed from
an independent hip-symmetry token.

\subsection{Autoregressive Generation Comparison}
\label{sec:generation_comparison}

The free-running comparison is conducted on the 310
geometry-grounded test records. YogaLLaVA V3 is geometry-conditioned
and task-specifically trained; LLaVA-1.5-7B and Qwen2.5-VL-7B are
evaluated zero-shot with image and question only. The comparison is
therefore not modality matched and does not isolate the causal
contribution of geometry conditioning, instruction tuning, or
continuous angle regression.

\paragraph{Model-specific slot-count coverage.}
Table~\ref{tab:generation_coverage} reports the number of records on
which each model emits the correct number of numeric slots. Numeric
errors in this table are calculated only over each model's own
successful subset and are therefore not directly comparable.

\begin{table}[t]
    \centering
    \caption{Model-specific free-generation slot-count coverage on 310
    geometry-grounded test records. Numeric metrics are calculated on
    different model-specific successful subsets.}
    \label{tab:generation_coverage}
    \resizebox{\columnwidth}{!}{
    \begin{tabular}{lrrrrr}
        \toprule
        \textbf{Model} &
        \textbf{Correct Records} &
        \textbf{Slot Match} &
        \textbf{Scored Slots} &
        \textbf{MAE $\downarrow$} &
        \textbf{RMSE $\downarrow$} \\
        \midrule
        YogaLLaVA V3 & 267/310 & 86.1\% & 365 & $25.14^\circ$ & -- \\
        LLaVA-1.5-7B & 205/310 & 66.1\% & 283 & $45.80^\circ$ & $56.30^\circ$ \\
        Qwen2.5-VL-7B & 211/310 & 68.1\% & 253 & $45.46^\circ$ & $58.21^\circ$ \\
        \bottomrule
    \end{tabular}}
\end{table}

YogaLLaVA V3 produces the correct number of values on 267 of 310
records, compared with 205 for LLaVA and 211 for Qwen. The
model-specific MAEs should not be read as a fair standalone accuracy
ranking because the models are scored on different records and slots.
Successful-subset filtering can favor a model that fails to produce
valid slot counts on more difficult examples.

For completeness, LLaVA and Qwen achieve within-$10^\circ$ rates of
15.55\% and 15.42\%, respectively, on their own successful subsets.
The corresponding value was not supplied for YogaLLaVA V3 and is not
inferred.

\paragraph{Common-success comparison.}
To compare numeric accuracy on aligned examples, we evaluate the 160
records on which all three models generate the correct number of
numeric slots.

\begin{table}[t]
    \centering
    \caption{Free-generation numeric accuracy on the common-success
    subset of 160 records for which all three models produced the
    correct number of numeric slots.}
    \label{tab:common_success}
    \begin{tabular}{lrr}
        \toprule
        \textbf{Model} &
        \textbf{Slot MAE $\downarrow$} &
        \textbf{Slot RMSE $\downarrow$} \\
        \midrule
        YogaLLaVA V3 & $17.65^\circ$ & $27.23^\circ$ \\
        LLaVA-1.5-7B & $42.80^\circ$ & $52.59^\circ$ \\
        Qwen2.5-VL-7B & $46.91^\circ$ & $61.12^\circ$ \\
        \bottomrule
    \end{tabular}
\end{table}

On this shared subset, YogaLLaVA V3 obtains lower slot-level error than
the two zero-shot baselines. Record-level macro MAE and bootstrap
intervals are shown in Table~\ref{tab:bootstrap_results}.

\begin{table}[t]
    \centering
    \caption{Record-level macro MAE with 95\% bootstrap confidence
    intervals on the 160-record common-success subset.}
    \label{tab:bootstrap_results}
    \begin{tabular}{lrr}
        \toprule
        \textbf{Model} &
        \textbf{Macro MAE $\downarrow$} &
        \textbf{95\% CI} \\
        \midrule
        YogaLLaVA V3 & $16.98^\circ$ & $[14.15,\,20.17]^\circ$ \\
        LLaVA-1.5-7B & $41.84^\circ$ & $[37.33,\,46.41]^\circ$ \\
        Qwen2.5-VL-7B & $47.63^\circ$ & $[41.59,\,53.87]^\circ$ \\
        \bottomrule
    \end{tabular}
\end{table}

The paired macro-MAE difference is $24.85^\circ$ for
LLaVA minus YogaLLaVA V3, with a 95\% confidence interval of
$[19.45,\,30.35]^\circ$. For Qwen minus YogaLLaVA V3, the difference is
$30.64^\circ$, with a 95\% confidence interval of
$[23.70,\,37.83]^\circ$.

These results show an advantage for the complete YogaLLaVA V3 system
on the aligned common-success subset. They do not establish superiority
under identical inputs: YogaLLaVA V3 receives structured geometry and
task-specific training, while the baselines receive only image and
question in a zero-shot setting. The comparison therefore measures the
combined effect of structured geometry, task-specific instruction
tuning, and continuous angle-slot regression.

\paragraph{Baseline language overlap.}
Table~\ref{tab:baseline_language_results} reports lexical-overlap
metrics for the two zero-shot baselines. No corresponding saved
YogaLLaVA V3 language-overlap result is claimed.

\begin{table}[t]
    \centering
    \caption{Lexical-overlap metrics for the zero-shot baselines under
autoregressive generation.}
    \label{tab:baseline_language_results}
    \begin{tabular}{lccc}
        \toprule
        \textbf{Model} &
        \textbf{BLEU-4 $\uparrow$} &
        \textbf{ROUGE-L $\uparrow$} &
        \textbf{Token F1 $\uparrow$} \\
        \midrule
        LLaVA-1.5-7B & 0.103 & 0.247 & 0.313 \\
        Qwen2.5-VL-7B & 0.091 & 0.236 & 0.318 \\
        \bottomrule
    \end{tabular}
\end{table}

These lexical metrics are not direct evidence of numeric or geometric
correctness. A response may overlap with the reference while containing
an inaccurate angle, anatomical side, comparison, or conclusion.

\subsection{Training Dynamics}
\label{sec:training_dynamics}

\begin{table}[t]
    \centering
    \caption{YogaLLaVA V3 training and teacher-forced validation
    performance. Epoch 10 is selected by minimum validation angle RMSE.}
    \label{tab:training_results}
    \resizebox{\columnwidth}{!}{
    \begin{tabular}{crrrrr}
        \toprule
        \textbf{Epoch} &
        \textbf{Train Loss} &
        \textbf{Val. Loss} &
        \textbf{Val. RMSE $\downarrow$} &
        \textbf{Val. MAE $\downarrow$} &
        \textbf{$R^2$ $\uparrow$} \\
        \midrule
        1  & 1.6154 & 1.1739 & $31.42^\circ$ & $21.40^\circ$ & 0.683 \\
        2  & 1.0768 & 1.1369 & $28.75^\circ$ & $20.61^\circ$ & 0.735 \\
        3  & 0.9607 & 1.1370 & $24.32^\circ$ & $16.59^\circ$ & 0.810 \\
        4  & 0.8407 & 1.1841 & $24.81^\circ$ & $15.79^\circ$ & 0.803 \\
        5  & 0.7323 & 1.2609 & $21.62^\circ$ & $14.25^\circ$ & 0.850 \\
        6  & 0.6481 & 1.3490 & $21.18^\circ$ & $13.72^\circ$ & 0.856 \\
        7  & 0.5976 & 1.3969 & $20.79^\circ$ & $13.03^\circ$ & 0.861 \\
        8  & 0.5624 & 1.4319 & $20.07^\circ$ & $13.02^\circ$ & 0.871 \\
        9  & 0.5405 & 1.4645 & $19.94^\circ$ & $12.74^\circ$ & 0.872 \\
        \textbf{10} & \textbf{0.5271} & 1.4821 &
        $\mathbf{19.83^\circ}$ & $\mathbf{12.63^\circ}$ & \textbf{0.874} \\
        11 & 0.5198 & 1.5038 & $19.93^\circ$ & $12.65^\circ$ & 0.873 \\
        12 & 0.5175 & 1.5082 & $20.03^\circ$ & $12.69^\circ$ & 0.871 \\
        \bottomrule
    \end{tabular}}
\end{table}

Training loss decreases steadily across all 12 epochs. Validation angle
RMSE also improves substantially and reaches its minimum at epoch 10.
In contrast, total validation loss is lowest around epochs 2--3 and
then increases. Because the total validation objective combines
language cross-entropy with weighted regression loss and is dominated
by the language term, this pattern is consistent with language-side
overfitting while numeric validation RMSE continues to improve. It is
not, by itself, definitive proof of the cause.

The selected epoch-10 checkpoint achieves validation RMSE
$19.8274^\circ$, MAE $12.6342^\circ$, and $R^2=0.8739$. The selection
policy optimizes teacher-forced numeric validation RMSE rather than
free-running generated-language quality.

\subsection{Interpretation and Scope}
\label{sec:result_interpretation}

YogaLLaVA V3 achieves a teacher-forced test MAE of $13.98^\circ$ and
RMSE of $21.82^\circ$ across 521 known angle slots, with 57.01\% of
predictions falling within $10^\circ$ of the reconstruction-derived
references. These values measure regression at known reference slot
locations and do not equal complete free-running inference performance.

During free generation, the model must additionally determine whether a
numeric slot is needed, emit the correct number of
\texttt{<ANGLE>} tokens, associate and order them correctly, generate
coherent prose, maintain consistency between numerical and qualitative
content, and terminate properly. The generation results are therefore
reported separately from teacher-forced regression.

For 518 of 521 test targets, the exact requested metric and value are
already present in the structured geometry input. The principal
teacher-forced result consequently supports structured-geometry routing,
question-conditioned selection, numeric reproduction, and
language-conditioned expression. It does not demonstrate independent
image-only recovery of physical joint angles.

The references are SAM 3D Body-derived pseudo-measurements and may be
affected by occlusion, viewpoint, foreshortening, loose clothing,
unusual inversions, left--right ambiguity, and monocular depth
uncertainty. The results should not be interpreted as clinical,
diagnostic, rehabilitation, motion-capture, or manually verified
biomechanical accuracy. Finally, because the test split contains unseen
images but shares pose categories with training, the experiments do not
establish generalization to unseen yoga poses.
